%%% SECTION 1: INTRODUCTION TO PHYSICAL AI ONCOLOGY TRIALS %%%

[This introduction should be approximately 8-10 pages and establish the full
context for why a National Platform for Physical AI Oncology Trials is needed.
The major theme is that Physical AI represents a transformative integration of
robotics and artificial intelligence into oncology clinical trials that will
enable more patients to be treated with more robots and fewer workers, while
shifting control towards the patient's side with expanded patient rights.]

\subsection{Background and Motivation}
\label{subsec:background}

[Build this subsection using information from
national-platform/research\_a/01\_research\_a\_part1.txt (the executive summary
section on AI in clinical trials) and
national-platform/research\_b/01\_research\_b\_part1.txt (the executive summary
on California and federal legal frameworks). Establish that oncology drug
development faces severe delays and costs, referencing Tufts CSDD data
\cite{tufts-delay} on the value of each day of delay. Reference existing FDA
clinical trials guidance \cite{fda-clinical-trials-guidance} and FDA oncology
guidance \cite{fda-oncology-guidance} to show the current landscape. Cite NCI
clinical trials information \cite{nci-clinical-trials} and ClinicalTrials.gov
\cite{clinicaltrials-gov} for current trial infrastructure context.]

[MAJOR THEME: Current oncology clinical trials are slow, expensive, and
resource-constrained. Physical AI offers a path to treat significantly more
patients at lower cost and higher quality. Reference the 24-hour simulation
evidence from Clinical Trial Site Complete Documentation \cite{site-docs}
showing 168 patients treated with 29 robots at 99.7 percent uptime.]

\subsection{Defining Physical AI in Oncology}
\label{subsec:defining-pai}

[Define Physical AI as the integration of embodied robotic systems with
advanced AI for direct patient care in clinical trial settings. Use definitions
from national-platform/21cfr312\_adapt/01\_preamble\_scope\_definitions.tex
(which contains 22 Physical AI definitions in \S 312.3) and
national-platform/21cfr50\_adapt/01\_preamble\_scope\_definitions\_consent.tex
(Physical AI definitions in Subpart A). Also reference the FDA AI glossary
\cite{fda-ai-glossary} for official FDA terminology. Reference Claude Opus 4.6
\cite{claude-opus}, GPT-5.4 \cite{gpt-5-4}, and Gemini \cite{gemini} as
state of the art AI systems that enable Physical AI development.]

\subsection{The Case for a National Platform}
\label{subsec:case-national}

[MAJOR THEME: Explain why this National Platform is more comprehensive and
meaningful compared to existing FDA approaches. Unlike FDA's piecemeal guidance
documents that address individual topics in isolation, this platform provides:
(a) three simultaneously adapted core regulatory standards (ICH E6(R3), 21 CFR
Part 50, 21 CFR Part 312) that build on existing well-established frameworks
and add credibility to Physical AI implementations; (b) complete site
establishment documentation with eleven regulatory documents; (c) validated
simulation evidence at both single-patient and 168-patient scales;
(d) quantitative robot readiness standards through PSL and USL; and
(e) national infrastructure for multi-site coordination through MCP servers
and federated learning.]

[Use national-platform/national\_mcp/chunk1\_preamble\_and\_introduction.tex
(the introduction section on fragmentation) to describe the current
fragmentation problem. Reference the Physical AI National MCP Servers
\cite{national-mcp-paper} and Physical AI Federated Learning \cite{fl-paper}
as solutions.]

[Reference NCI modernizing clinical trials \cite{nci-modernizing} and FDA
adaptive designs guidance \cite{fda-adaptive-designs} to show that even
established institutions recognize the need for modernization, but none have
proposed a solution as comprehensive as this National Platform.]

\subsection{Scope and Organization of This Document}
\label{subsec:scope}

[Provide a roadmap of the paper's sixteen sections. Briefly describe what each
section covers and how sections build upon each other. Note that this document
draws from the complete body of work listed in Source Documents and Their
Significance, including three adapted regulatory standards (Adaption: ICH
Harmonised Guideline \cite{ich-adapt}, Physical AI Adaption: 21 CFR Part 50
\cite{cfr50-adapt}, Physical AI Adaption: 21 CFR Part 312 \cite{cfr312-adapt}),
two standards frameworks (Physical AI Unification Standard Level
\cite{usl-standard} and Clinical Trial Site Complete Documentation
\cite{site-docs}), two patient-centered documents (A Cancer Patient's Journey,
Physical AI \cite{patient-journey-paper} and Patient Instructions: Physical AI
Trials \cite{patient-instructions-paper}), two infrastructure papers (Physical
AI National MCP Servers \cite{national-mcp-paper} and Physical AI Federated
Learning \cite{fl-paper}), and three code repositories
\cite{main-repo, mcp-repo, fl-repo}.]

\subsection{Contributions}
\label{subsec:contributions}

[List the key contributions of this National Platform as a numbered or bulleted
list. Include: (1) first comprehensive end-to-end regulatory framework for
Physical AI oncology trials; (2) adapted ICH/CFR standards with Physical AI
provisions; (3) quantitative PSL and USL scoring frameworks; (4) validated
simulation evidence at scale; (5) complete site establishment documentation;
(6) national MCP server architecture; (7) federated learning for multi-site
privacy-preserving coordination; (8) patient-centered trial design with
expanded rights; (9) economic analysis demonstrating cost and time savings;
(10) open-source implementation across three repositories.]
